% ============================================================================
% УВОД, без номер
% ============================================================================
\section*{УВОД}
\addcontentsline{toc}{section}{УВОД}

Видеозаписът е едно от най-често срещаните веществени доказателства в съвременното
разследване. Източниците са много: камери за видеонаблюдение, видеорегистратори в
автомобили, нагръдни камери, мобилни телефони. Общото между тях е, че записът
съществува като файл върху цифров носител. Затова преди да започне анализът на
съдържанието (лица, регистрационни номера, последователност на събития), експертът
трябва да разполага с файл, който може да бъде прочетен и възпроизведен.

На практика това често не е така. Файлът може да е изтрит от потребителя, носителят
може да е форматиран, записът може да е прекъснат от спиране на захранването, а
части от файла може да са презаписани с нови данни. В тези случаи стандартните
програми за възпроизвеждане отказват да отворят файла или показват само част от него.
Възстановяването на такива файлове е отделна задача на цифровата криминалистика, която
стои преди същинския видеоанализ.

\term{Целта на курсовата работа} е да опише теоретичната основа на възстановяването на
изтрити и повредени видеофайлове и да оцени доказателствената стойност на резултата.
За постигането ѝ са поставени следните задачи:

\begin{enumerate}[leftmargin=2cm]
  \item да се опише структурата на видеофайл във формат MP4 и на видеопотока H.264,
        доколкото тя определя какво може да бъде възстановено;
  \item да се систематизират причините за загуба и повреда на видеофайлове;
  \item да се сравнят основните методи за възстановяване: по метаданни на файловата
        система, чрез изрязване по сигнатури (\term{file carving}), за фрагментирани
        файлове, на ниво кадри и чрез ремонт с еталонен файл;
  \item да се опишат изискванията към процедурата, така че възстановеният запис да
        запази доказателствената си стойност.
\end{enumerate}

Работата е теоретична. Тя не съдържа програмен код и собствени измервания. Всички
технически твърдения се позовават на стандартите и на публикуваните източници,
посочени в списъка с литература.
